\chapter{Formal Declaration Crosswalk and Mechanization Status}
\label{app:formal-crosswalk}
\small
\sloppy
\hbadness=10000

\section{Purpose of the crosswalk}

This appendix binds the mathematical exposition to stable declaration identifiers,
exact Lean names, source modules, and statement-level proof dossiers. The identifiers
are editorial handles; the Lean declarations and their proof terms remain authoritative.
A range in the tables means that every identifier in the inclusive range is represented
in the theorem manifest. It does not authorize a stronger aggregate statement than the
individual declarations prove.

\section{Finite-frame sequence}

The definitions and ten theorems in Chapter~\ref{ch:finite-operator-foundations} are
implemented in
\modulepath{ProjectKoios/BandGauge/Alignment/FiniteFrame.lean}. Their statement dossiers
are stored under \modulepath{proofs/BandGauge/Alignment/}.

\begin{description}
  \item[\decl{BG-ALIGNMENT-THM-001}.]
    \path{ProjectKoios.BandGauge.Alignment.frameProjector_isHermitian};
    dossier \modulepath{BG-ALIGNMENT-THM-001-projector-hermitian.tex}.
    This is Theorem~\ref{thm:bg-alignment-001}; it has no orthonormality premise.

  \item[\decl{BG-ALIGNMENT-THM-002}.]
    \path{ProjectKoios.BandGauge.Alignment.frameProjector_mul_self};
    dossier \modulepath{BG-ALIGNMENT-THM-002-projector-idempotent.tex}.
    This is Theorem~\ref{thm:bg-alignment-002}; exact frame orthonormality is required.

  \item[\decl{BG-ALIGNMENT-THM-003}.]
    \path{ProjectKoios.BandGauge.Alignment.rightUnitaryFrame_isOrthonormal};
    dossier \modulepath{BG-ALIGNMENT-THM-003-right-unitary-orthonormal.tex}.
    This is Theorem~\ref{thm:bg-alignment-003}.

  \item[\decl{BG-ALIGNMENT-THM-004}.]
    \path{ProjectKoios.BandGauge.Alignment.frameProjector_rightUnitaryFrame};
    dossier \modulepath{BG-ALIGNMENT-THM-004-projector-right-unitary.tex}.
    This is Theorem~\ref{thm:bg-alignment-004}; it does not require orthonormality.

  \item[\decl{BG-ALIGNMENT-THM-005}.]
    \path{ProjectKoios.BandGauge.Alignment.compressOperator_rightUnitaryFrame};
    dossier \modulepath{BG-ALIGNMENT-THM-005-compression-covariance.tex}.
    This is Theorem~\ref{thm:bg-alignment-005}.

  \item[\decl{BG-ALIGNMENT-THM-006}.]
    \path{ProjectKoios.BandGauge.Alignment.compressOperator_isHermitian};
    dossier \modulepath{BG-ALIGNMENT-THM-006-compression-hermitian.tex}.
    This is Theorem~\ref{thm:bg-alignment-006}; frame orthonormality is not required.

  \item[\decl{BG-ALIGNMENT-THM-007}.]
    \path{ProjectKoios.BandGauge.Alignment.reconstructCompressedOperator_eq_projector_sandwich};
    dossier \modulepath{BG-ALIGNMENT-THM-007-reconstruction-projector-sandwich.tex}.
    This is Theorem~\ref{thm:bg-alignment-007}.

  \item[\decl{BG-ALIGNMENT-THM-008}.]
    \path{ProjectKoios.BandGauge.Alignment.reconstructCompressedOperator_rightUnitaryFrame};
    dossier \modulepath{BG-ALIGNMENT-THM-008-reconstruction-right-unitary.tex}.
    This is Theorem~\ref{thm:bg-alignment-008}; it does not require orthonormality.

  \item[\decl{BG-ALIGNMENT-THM-009}.]
    \path{ProjectKoios.BandGauge.Alignment.frameProjector_mul_frame};
    dossier \modulepath{BG-ALIGNMENT-THM-009-projector-fixes-frame.tex}.
    This is Theorem~\ref{thm:bg-alignment-009}.

  \item[\decl{BG-ALIGNMENT-THM-010}.]
    \path{ProjectKoios.BandGauge.Alignment.compressOperator_frameProjector};
    dossier \modulepath{BG-ALIGNMENT-THM-010-compressed-projector-identity.tex}.
    This is Theorem~\ref{thm:bg-alignment-010}.
\end{description}

The definitions used by the sequence are
\decl{IsOrthonormalFrame}, \decl{frameProjector}, \decl{rightUnitaryFrame},
\decl{compressOperator}, and \decl{reconstructCompressedOperator}. Exact
orthonormality is a hypothesis, not a conclusion about a runtime frame.

\section{Bloch alignment and alignment families}

\begin{table}[ht]
\centering
\small
\begin{tabularx}{\textwidth}{@{}Z Y Y@{}}
\toprule
Declaration range & Source module & Mathematical role \\
\midrule
\decl{BG-BLOCH-THM-012--020}
  & \modulepath{ProjectKoios/BandGauge/Bloch/AlignedDifference.lean}
  & pointwise difference identities, shared-sewing closure, sign reversal, recovery,
    zero criteria, and aligned-difference sewing \\
\decl{BG-ALIGNMENT-DEF-007--016}
  & \modulepath{ProjectKoios/BandGauge/Alignment/Family.lean}
  & singleton, identity, global-unitary, and complete pointwise alignment families and
    their feasibility predicates \\
\decl{BG-ALIGNMENT-THM-011--019}
  & \modulepath{ProjectKoios/BandGauge/Alignment/Family.lean}
  & family membership, forward inclusions, and feasibility transfer \\
\decl{BG-ALIGNMENT-THM-020--029}
  & \modulepath{ProjectKoios/BandGauge/Alignment/FamilySeparation.lean}
  & strict family separation and exact reverse-feasibility counterexamples \\
\bottomrule
\end{tabularx}
\caption{Aligned-difference and alignment-family implementation crosswalk.}
\label{tab:alignment-family-implementation-crosswalk}
\end{table}

The corresponding dossiers are under \modulepath{proofs/BandGauge/Bloch/} and
\modulepath{proofs/BandGauge/Alignment/}. The strictness declarations use supplied point
and unitary inequalities. Their constructed acceptance predicates establish logical
separation of decision spaces; they do not certify the outcome of a scientific
application.

\section{Three-channel residual calculus}

The source-independent finite residual algebra used in
Theorem~\ref{thm:three-channel-residual-calculus} is implemented in
\modulepath{ProjectKoios/BandGauge/Hilbert/ResidualChannels.lean}.

\begin{table}[ht]
\centering
\small
\begin{tabularx}{\textwidth}{@{}Z Y Y@{}}
\toprule
Declaration range & Formal objects & Mathematical role \\
\midrule
\decl{BG-HILBERT-DEF-013--016}
  & \decl{FiniteResidualField}, \decl{interpolationResidual},
    \decl{truncationResidual}, \decl{totalResidual}
  & the common finite field type and the three signed residual channels \\
\decl{BG-HILBERT-DEF-017--018}
  & \decl{finiteMaximumNorm}, \decl{finiteRootMeanSquareNorm}
  & the finite maximum and root-mean-square summaries \\
\decl{BG-HILBERT-THM-016--017}
  & maximum-as-supremum and explicit RMS square-sum formula
  & normalization and interpretation of the two aggregate norms \\
\decl{BG-HILBERT-THM-018--021}
  & exact decomposition and pointwise, maximum, and RMS bounds
  & the three-channel identity and its norm-triangle consequences \\
\bottomrule
\end{tabularx}
\caption{Finite three-channel residual implementation crosswalk.}
\label{tab:three-channel-implementation-crosswalk}
\end{table}

The four principal declarations are:

\begin{description}
  \item[\decl{BG-HILBERT-THM-018}.]
    \path{ProjectKoios.BandGauge.Hilbert.totalResidual_eq_add};
    dossier \modulepath{BG-HILBERT-THM-018-residual-decomposition.tex}.
  \item[\decl{BG-HILBERT-THM-019}.]
    \path{ProjectKoios.BandGauge.Hilbert.norm_totalResidual_le_add};
    dossier \modulepath{BG-HILBERT-THM-019-pointwise-triangle.tex}.
  \item[\decl{BG-HILBERT-THM-020}.]
    \path{ProjectKoios.BandGauge.Hilbert.finiteMaximumNorm_totalResidual_le_add};
    dossier \modulepath{BG-HILBERT-THM-020-maximum-triangle.tex}.
  \item[\decl{BG-HILBERT-THM-021}.]
    \path{ProjectKoios.BandGauge.Hilbert.finiteRootMeanSquareNorm_totalResidual_le_add};
    dossier \modulepath{BG-HILBERT-THM-021-rms-triangle.tex}.
\end{description}

The dossiers are stored under \modulepath{proofs/BandGauge/Hilbert/}.  The exact identity
requires one common intermediate field.  The aggregate theorems are inequalities and do
not assert additivity of reported scalar errors.  The formal RMS definition is totalized
on an empty inventory; the scientific formulas in the main text use the ordinary
nonempty finite-inventory interpretation.

\section{Fourier and gauge families}

\begin{table}[ht]
\centering
\small
\begin{tabularx}{\textwidth}{@{}Z Y Y@{}}
\toprule
Declaration range & Source module & Mathematical role \\
\midrule
\decl{BG-FOURIER-THM-033--047}
  & \modulepath{ProjectKoios/BandGauge/Fourier/UnitaryEquivariance.lean}
  & fixed-unitary preservation and commutation with the finite Fourier and truncation
    constructions \\
\decl{BG-FOURIER-THM-048--054}
  & \modulepath{ProjectKoios/BandGauge/Fourier/PointwiseGauge.lean}
  & mode shift, scalar-conjugation boundary, coefficient change, and pointwise
    Frobenius-power preservation \\
\decl{BG-FOURIER-THM-055--066}
  & \modulepath{ProjectKoios/BandGauge/Fourier/GaugeTruncationCounterexample.lean}
  & exact two-point rank-two gauge--truncation counterexample \\
\decl{BG-FOURIER-THM-067--075}
  & \modulepath{ProjectKoios/BandGauge/Fourier/ResidualChannels.lean}
  & deliberate-truncation residual, omitted power, root-mean-square identities, and
    pointwise and finite-maximum bounds \\
\bottomrule
\end{tabularx}
\caption{Fourier and gauge implementation crosswalk.}
\label{tab:fourier-gauge-implementation-crosswalk}
\end{table}

The key declarations for Theorem~\ref{thm:two-point-gauge-truncation} are:

\begin{description}
  \item[\decl{BG-FOURIER-THM-062}.]
    \path{ProjectKoios.BandGauge.Fourier.pointwiseSwapConjugation_odd_coefficient};
    dossier \modulepath{BG-FOURIER-THM-062-gauged-odd-coefficient.tex}.
  \item[\decl{BG-FOURIER-THM-064}.]
    \path{ProjectKoios.BandGauge.Fourier.pointwiseSwapConjugation_odd_coefficient_ne_zero};
    dossier \modulepath{BG-FOURIER-THM-064-gauged-odd-mode-nonzero.tex}.
  \item[\decl{BG-FOURIER-THM-066}.]
    \path{ProjectKoios.BandGauge.Fourier.pointwiseSwapConjugation_zeroModeFourierTruncation_ne};
    dossier \modulepath{BG-FOURIER-THM-066-gauge-truncation-noncommutation.tex}.
\end{description}

The declaration sequence for Theorem~\ref{thm:finite-fourier-tail-bounds} is:

\begin{description}
  \item[\decl{BG-FOURIER-THM-067}.]
    \path{ProjectKoios.BandGauge.Fourier.truncationResidual_fourierReconstruction_eq}.
  \item[\decl{BG-FOURIER-THM-068}.]
    \path{ProjectKoios.BandGauge.Fourier.finiteFourierPower_truncationResidual_eq_omitted}.
  \item[\decl{BG-FOURIER-THM-069}.]
    \path{ProjectKoios.BandGauge.Fourier.finiteMatrixFrobeniusPower_truncationResidual_eq_omitted}.
  \item[\decl{BG-FOURIER-THM-070}.]
    \path{ProjectKoios.BandGauge.Fourier.finiteRootMeanSquareNorm_truncationResidual_eq_sqrt_omitted}.
  \item[\decl{BG-FOURIER-THM-071}.]
    \path{ProjectKoios.BandGauge.Fourier.finiteRootMeanSquareNorm_truncationResidual_sq_eq_omitted}.
  \item[\decl{BG-FOURIER-THM-072}.]
    \path{ProjectKoios.BandGauge.Fourier.finiteMatrixFrobeniusRootMeanSquare_truncationResidual_eq_sqrt_omitted}.
  \item[\decl{BG-FOURIER-THM-073}.]
    \path{ProjectKoios.BandGauge.Fourier.finiteMatrixFrobeniusRootMeanSquare_truncationResidual_sq_eq_omitted}.
  \item[\decl{BG-FOURIER-THM-074}.]
    \path{ProjectKoios.BandGauge.Fourier.norm_truncationResidual_fourierReconstruction_le_omitted}.
  \item[\decl{BG-FOURIER-THM-075}.]
    \path{ProjectKoios.BandGauge.Fourier.finiteMaximumNorm_truncationResidual_le_omitted}.
\end{description}

The dossiers follow the stable filenames
\modulepath{proofs/BandGauge/Fourier/BG-FOURIER-THM-067-truncation-channel.tex}
through
\modulepath{proofs/BandGauge/Fourier/BG-FOURIER-THM-075-finite-maximum-tail.tex}.
The two root-mean-square families require the explicit nonempty-domain premise recorded
in their Lean types. The maximum-norm result uses the ambient norm and is not an
entrywise matrix-Frobenius specialization.

\section{Reduction, certificates, and finite decisions}

\begin{table}[ht]
\centering
\small
\begin{tabularx}{\textwidth}{@{}Z Y Y@{}}
\toprule
Declaration range & Principal source & Scope represented in Chapter~\ref{ch:reduction-decision} \\
\midrule
\decl{BG-THM-001--037}
  & \modulepath{ProjectKoios/BandGauge/Geometry/}
  & common witnesses, compact attainment, zero/positive separation converses, rank-two
    admissible regions, and exact set distances \\
\decl{BG-HILBERT-THM-001--015}
  & \modulepath{ProjectKoios/BandGauge/Hilbert/}
  & ball geometry, weighted Frobenius algebra, and spectral/operator residual bridges \\
\decl{BG-FOUNDATION-THM-008--014}
  & \modulepath{ProjectKoios/BandGauge/Foundations/RankTwoSpectralPerturbation.lean}
  & canonical Hermitian rank-two spectra and Frobenius perturbation bounds \\
\decl{BG-CERT-THM-001--009}, \decl{BG-CERT-THM-013--016}
  & \modulepath{ProjectKoios/BandGauge/Certificates/}
  & exact, enclosure, and box separation certificates \\
\decl{BG-CERT-THM-010--012}
  & \modulepath{ProjectKoios/BandGauge/Certificates/FinitePareto.lean}
  & exact finite dominance, feasibility, and efficient-mask checks \\
\decl{BG-CERT-THM-017--025}
  & \modulepath{ProjectKoios/BandGauge/Certificates/IntervalDecision.lean}
  & accepted, rejected, and unresolved interval decisions and their soundness \\
\decl{BG-DECISION-THM-014--016}
  & \modulepath{ProjectKoios/BandGauge/Decision/FiniteCutoff.lean}
  & first passing and first joint-passing semantics on a finite inventory \\
\decl{BG-DECISION-DEF-012--027}
  & \modulepath{ProjectKoios/BandGauge/Decision/FiniteCompatibility.lean}
  & finite model/candidate inventories, infimum and witness compatibility, mandatory
    channels, conformant embeddings, metric separation, and first-compatible selection \\
\decl{BG-DECISION-THM-017--038}
  & \modulepath{ProjectKoios/BandGauge/Decision/FiniteCompatibility.lean}
  & infimum/attainment bridges, common-candidate and common-alignment semantics,
    compatibility transfer, finite-set distance attainment and zero/positive converses,
    and finite first-compatible theorems \\
\decl{QM-RESOLVENT-THM-001--014}
  & \modulepath{ProjectKoios/QuantumMechanics/FiniteDimensional/RankOneResolvent.lean}
  & exact rank-one resolvent, determinant, secular, kernel, and Hermiticity identities \\
\decl{SS-GAUGE-THM-001--023}
  & \modulepath{ProjectKoios/SolidState/FiniteLattice/}
  & finite gauge loops, determinant invariance, projector products, and Bargmann
    identities \\
\decl{SS-ROUTE-THM-001--010}
  & \modulepath{ProjectKoios/SolidState/FiniteLattice/RouteCommutativity.lean}
  & finite pullback, spin, folding, and defect-extraction route algebra \\
\bottomrule
\end{tabularx}
\caption{Declaration families supporting the reduction and decision narrative.}
\label{tab:reduction-decision-implementation-crosswalk}
\end{table}
\FloatBarrier

These declarations now package the generic finite spectral--operator compatibility
logic, including alignment-infimum and stronger common-witness semantics, explicit
model embeddings, finite central admissible sets, attained separation under nonemptiness,
zero/positive distance converses, and first-compatible selection.  No declaration in the
pinned snapshot instantiates the schema with the complete ICMSEP model hierarchy,
application-specific losses, tolerances, supports, channel contracts, or retained
certificate.  Chapter~\ref{ch:mechanization-status} records that proposition-specific
assembly as open.

\section{Mechanization status at the pinned snapshot}

The theorem manifest is the authoritative declaration inventory. At the pinned snapshot,
the repository audit covers 645 proved theorem or lemma rows and 649 unique named axiom
reports. The repository-owned Lean sources contain no \texttt{sorry}, \texttt{admit}, or
repository-owned \texttt{axiom} declaration. Individual reports may include the standard
foundational dependencies \texttt{propext}, \texttt{Classical.choice}, and
\texttt{Quot.sound}; the exact report for each declaration is recorded in the theorem
manifest and imported by \modulepath{ProjectKoios/Audit.lean}.

Checked reusable families include finite Fourier transformation and truncation, finite
frames and alignment families, weighted Frobenius and admissible-set geometry, exact and
interval certificates, finite first-passing decisions, finite lattice folding and route
algebra, rank-one resolvents, rank-two perturbation, finite gauge loops and Bargmann
identities, and finite quantum-state, measurement, and partial-operator foundations.
Membership in this inventory means that the exact Lean proposition is proved. It does
not imply runtime conformance, continuum validity, physical applicability, novelty, or
publication status.
